\section{Discussion}\label{sec:discussion}

\paragraph{Relation to other model families.}
Stochastic actor-oriented models \citep{snijders2001,snijders2010,snijders2017} describe network
change by actors who reconsider one outgoing tie at a time according to an objective function, and
are estimated from panel data by the method of moments or by likelihood-based methods with data
augmentation \citep{koskinen2007}. Relational event models \citep{butts2008} and dynamic network actor
models \citep{stadtfeld2017} model sequences of time-stamped interactions with rates that depend on
statistics of the past. \tricl\ is tie-oriented rather than actor-oriented: each potential link has
its own establishment and termination rates, which depend on the current two-path structure. It
shares with relational event models the exact likelihood of a fully observed event sequence, and
with the adaptive-network literature \citep{gross2008} the coupling of states and ties, which it
expresses within the same rule language through hub entities. Its distinctive features are typed
multi-relational networks with symmetric and inverse relationship types, models as configuration
files, immediate events for state logic, and the derivation of simulation, likelihood and mean-field
equations from the same file. The interval output of a run is a temporal network
\citep{holme2012} in the edge-list form that temporal-network software reads, and a mapping to the
Resource Description Framework (entity types as classes, relationship types as properties,
symmetric and inverse types as their OWL counterparts, link intervals as annotated triples) lets
\tricl\ import typed graphs from knowledge bases and export trajectories to them.

\paragraph{Limitations.}
Estimation requires a complete event sequence, including the timing of events. Most network data are
snapshots or panels; for them the likelihood is an integral over unobserved paths, which needs
simulation-based methods such as data augmentation or the method of simulated moments, both of which
the fast simulator makes feasible but which are not implemented yet. The rule language covers
two-path dependence only, and the coefficients do not depend on time; rules on single links of one
endpoint (degree dependence) are part of the syntax but not yet supported. The mean-field closure
ignores correlations between links (Section~\ref{sec:macro}). In the application, the yearly
resolution and the dyadic coding of multilateral treaties make the formation effect partly a
property of treaties rather than of pairs.

\paragraph{Outlook.}
Next steps are higher-order (pair and triple) closures generated in the same way, estimation from
snapshot data, rules on single links, and applications with event-level data, e.g.\ online
follow networks and scientific collaboration, where the typed relationships (co-authorship, citation,
shared affiliation) are central.
